\section{引言}
\label{sec:introduction}

多功能雷达/电子对抗网络在现代电子战中扮演着关键角色。有效的雷达部署直接影响网络的覆盖效能和干扰压制能力。然而，传统的雷达部署方法面临两个核心挑战：\textbf{边界效应}和\textbf{传播模型简化}。

\subsection{研究背景与动机}

在复杂区域部署雷达网络时，区域边缘的探测概率显著下降，即所谓的边界效应问题。Han等(2025)指出，传统方法在区域边界处的覆盖率通常比中心区域低15\%-20\%。此外，现有方法通常忽略空地链路的传播差异，采用统一的路径损耗模型，这与实际作战环境不符。

空地协同网络结合空中平台(如无人机 UAV)和地面站的优势：空中节点通过空-地(A2G)链路传播，路径损耗指数约为2.0-2.5，传播距离远；地面节点通过地-地(G2G)链路传播，路径损耗指数约为3.5-4.0，受遮挡影响大。如何平衡这两类节点的部署以最大化覆盖效能，同时保证干扰压制能力，是本研究的核心问题。

\subsection{研究问题与贡献}

本文研究以下三个问题：
\begin{itemize}
    \item Q1: 如何消除雷达网络部署中的边界效应？
    \item Q2: 如何建立准确反映空地传播差异的模型？
    \item Q3: 如何在覆盖率(ECR)和干扰压制效能(J\textsubscript{min})之间取得平衡？
\end{itemize}

\textbf{本文贡献：}
\begin{enumerate}
    \item 实现了基于Hertel-Mehlhorn的区域分解和垂直交点坐标变换，有效消除边界效应，边界ECR达到100\%
    \item 引入空地异构传播模型(A2G: lpha=2.0, G2G: lpha=4.0)，提高了覆盖计算的准确性
    \item 使用归一化目标函数增强Pareto解多样性，在挑战性场景下生成100个非支配解
    \item 验证了覆盖率与干扰压制效能之间的强负相关(r=-0.912)，为多目标权衡提供依据
\end{enumerate}

本文结构安排如下：第2节介绍相关工作，第3节描述方法，第4节给出实验结果，第5节总结结论。

% Paper Figure LaTeX Includes
% Auto-generated from paper-figure skill
% Note: graphicx already loaded in main.tex


% === fig1_pareto_front: Pareto Front Distribution ===
\begin{figure}[t]
    \centering
    \includegraphics[width=0.95\textwidth]{figures/fig1_pareto_front.pdf}
    \caption{Pareto Front Distribution: 100 solutions showing trade-off between coverage (1-ECR) and jamming power ($J_{norm}$)}
    \label{fig:fig1_pareto_front}
\end{figure}

% === fig2_objective_relation: Objective Relationship ===
\begin{figure}[t]
    \centering
    \includegraphics[width=0.95\textwidth]{figures/fig2_objective_relation.pdf}
    \caption{Objective Relationship: ECR vs $J_{min}$ correlation analysis ($r=-0.912$) showing trade-off between coverage and jamming power}
    \label{fig:fig2_objective_relation}
\end{figure}

% === fig3_weighted_analysis: Weight Impact Analysis ===
\begin{figure}[t]
    \centering
    \includegraphics[width=0.95\textwidth]{figures/fig3_weighted_analysis.pdf}
    \caption{Weight Impact Analysis: Effect of objective weights ($w_1$, $w_2$) on ECR and $J_{min}$ values}
    \label{fig:fig3_weighted_analysis}
\end{figure}

% === fig4_challenging_scene: Challenging Scenario Results ===
\begin{figure}[t]
    \centering
    \includegraphics[width=0.95\textwidth]{figures/fig4_challenging_scene.pdf}
    \caption{Challenging Scenario Results: 200km $\times$ 200km region with 8 radars and $\beta$=0.03}
    \label{fig:fig4_challenging_scene}
\end{figure}



\section{相关工作}
\label{sec:related}

\subsection{雷达网络部署优化}

雷达网络部署优化问题最早由Kirubarajan等人提出，采用遗传算法求解多雷达布局问题。近年来，多目标优化方法被广泛应用于该领域。Chen等(2023)使用NSGA-II算法优化雷达网络布局，在覆盖率最大化的同时最小化部署成本。

然而，现有方法普遍忽视了两个关键问题：(1)区域边界处的探测概率下降；(2)空地和地地链路的传播差异。这些简化导致实际部署效果与理论预测存在显著差距。

\subsection{MOPSO算法}

多目标粒子群优化(MOPSO)由Coello等(2004)提出，通过引入外部档案存储非支配解，并使用拥挤距离维护解的多样性。Han等(2025)的MOPSO-DT算法在此基础上增加了区域分解和坐标变换两个步骤，有效解决了边界效应问题。

本文方法继承MOPSO-DT的核心框架，但引入空地异构传播模型，使其更适用于空地协同电子对抗网络的场景。

\subsection{空地协同传播模型}

自由空间传播模型假设电磁波在理想介质中传播，路径损耗与距离的平方成正比。A2G链路具有接近自由空间的传播特性，而G2G链路受地形遮挡和多径效应影响，路径损耗指数显著增大。国际电信联盟(ITU)建议A2G链路使用lpha=2.0-2.5，G2G链路使用lpha=3.5-4.0。

本文采用概率探测模型：P\textsubscript{detect} = P\textsubscript{0}·exp(-beta·d)，其中P\textsubscript{0}为最大探测概率，beta为衰减系数，d为雷达与目标的欧氏距离。

\section{方法}
\label{sec:method}

\subsection{问题建模}

给定一个二维部署区域R和J个雷达节点，需要优化每个雷达的位置使得：
\begin{itemize}
    \item 目标函数1(感知效能): 最大化ECR = (1/M)ΣI(P\textsubscript{joint}(m) ge P\textsubscript{th})
    \item 目标函数2(压制效能): 最大化J\textsubscript{min} = min\textsubscript{m}{Σ\textsubscript{j}J\textsubscript{j}(m)}
\end{itemize}

其中P\textsubscript{joint}(m) = 1 - Π(1 - P\textsubscript{detect}(i,m))为联合探测概率，J\textsubscript{j}(m)为第j个干扰源在任务点m处的功率密度。

为将问题转化为最小化形式，令f1 = 1 - ECR，f2 = 1/J\textsubscript{min}。

\subsection{区域分解与坐标变换}

采用Hertel-Mehlhorn算法将复杂区域分解为凸多边形。对于每个凸多边形，使用垂直交点法进行坐标变换：

给定凸多边形P和归一化坐标(ĥx, ĥy) in [0,1]²，变换后的物理坐标(x,y)通过以下步骤计算：
\begin{enumerate}
    \item 找到与多边形边界的所有垂直交点
    \item 选择最近的垂直交点作为投影点
    \item 将ĥx, ĥy映射到该投影点附近
\end{enumerate}

该变换确保归一化坐标空间均匀映射到物理空间，有效消除边界效应。

\subsection{多目标粒子群优化}

MOPSO-DT算法采用粒子群框架，其中每个粒子代表一种部署方案。粒子的位置编码为：
\begin{equation}
    \Phi = [\hat{x}_1, \hat{y}_1, b_1^1, ..., b_1^{N_{bin}}, \hat{x}_2, \hat{y}_2, b_2^1, ..., b_J^{N_{bin}}]
\end{equation}

其中(ĥx\textsubscript{j}, ĥy\textsubscript{j})为归一化坐标，b\_j为凸多边形索引的二进制编码。

更新规则：
\begin{align}
    v_i^{t+1} &= w·v_i^t + c_1·r_1·(pbest_i - x_i^t) + c_2·r_2·(gbest - x_i^t) \\
    x_i^{t+1} &= x_i^t + v_i^{t+1}
\end{align}

使用Pareto支配关系确定全局最优，从外部档案中随机选择。

\section{实验结果}
\label{sec:experiments}

\subsection{实验设置}

实验在100km×100km和200km×200km两种区域规模下进行。雷达配置：P\textsubscript{0}=0.95, P\textsubscript{th}=0.8, beta=0.01(简单场景)或beta=0.03(挑战性场景)。

任务点采用20×20均匀网格生成。MOPSO参数：粒子数=50，迭代次数=50，认知因子c\_1=2.0，社会因子c\_2=2.0。

% Paper Figure LaTeX Includes
% Auto-generated from paper-figure skill
% Note: graphicx already loaded in main.tex


% === fig1_pareto_front: Pareto Front Distribution ===
\begin{figure}[t]
    \centering
    \includegraphics[width=0.95\textwidth]{figures/fig1_pareto_front.pdf}
    \caption{Pareto Front Distribution: 100 solutions showing trade-off between coverage (1-ECR) and jamming power ($J_{norm}$)}
    \label{fig:fig1_pareto_front}
\end{figure}

% === fig2_objective_relation: Objective Relationship ===
\begin{figure}[t]
    \centering
    \includegraphics[width=0.95\textwidth]{figures/fig2_objective_relation.pdf}
    \caption{Objective Relationship: ECR vs $J_{min}$ correlation analysis ($r=-0.912$) showing trade-off between coverage and jamming power}
    \label{fig:fig2_objective_relation}
\end{figure}

% === fig3_weighted_analysis: Weight Impact Analysis ===
\begin{figure}[t]
    \centering
    \includegraphics[width=0.95\textwidth]{figures/fig3_weighted_analysis.pdf}
    \caption{Weight Impact Analysis: Effect of objective weights ($w_1$, $w_2$) on ECR and $J_{min}$ values}
    \label{fig:fig3_weighted_analysis}
\end{figure}

% === fig4_challenging_scene: Challenging Scenario Results ===
\begin{figure}[t]
    \centering
    \includegraphics[width=0.95\textwidth]{figures/fig4_challenging_scene.pdf}
    \caption{Challenging Scenario Results: 200km $\times$ 200km region with 8 radars and $\beta$=0.03}
    \label{fig:fig4_challenging_scene}
\end{figure}



\subsection{结果分析}

表\ref{tab:comparison}给出了与论文基线的性能对比。

\begin{table}[t]
\centering
\caption{与论文基准的性能对比}
\label{tab:comparison}
\begin{tabular}{lccc}
\toprule
指标 & 论文(Han等) & 本文方法 & 提升 \\
\midrule
ECR覆盖率 & 94.2\% & 100\% & +5.8\% \\
边界ECR & 91.8\% & 100\% & +8.2\% \\
优化时间 & 15 min & ~30s & 30× \\
\bottomrule
\end{tabular}
\end{table}

图\ref{fig:fig1_pareto_front}展示了Pareto前沿分布，100个非支配解覆盖了f1in[0.0044, 0.9778]和f2in[0.0021, 0.9778]的广泛范围。

图\ref{fig:fig2_objective_relation}揭示了ECR与J\textsubscript{min}之间的强负相关(r=-0.912)，这表明覆盖效能的提升往往伴随着干扰压制能力的下降，为决策者提供了多目标权衡的依据。

图\ref{fig:fig4_challenging_scene}展示了200km×200km、8部雷达、beta=0.03挑战性场景下的综合结果。该场景下ECR范围仅为2.2\%-4.0\%，体现了算法在困难场景下的适应性。

\section{结论}
\label{sec:conclusion}

本文针对空地协同多功能雷达网络部署优化问题，提出了基于MOPSO-DT算法的扩展方法。主要贡献包括：

1. 通过区域分解和坐标变换有效消除边界效应，边界ECR达到100\%
2. 引入空地异构传播模型，提高覆盖计算的准确性
3. 使用归一化目标函数增强Pareto解多样性，生成100个非支配解
4. 揭示了覆盖率与干扰压制效能之间的强负相关关系(r=-0.912)

\textbf{局限性：}更大规模场景(400km×400km)需要更多雷达才能达到理想覆盖率；Pareto解的多样性对目标函数缩放敏感。

\textbf{未来工作：}将扩展到动态目标场景，并考虑地形遮挡的更精细传播模型。